% ---------------------------------------------------------------------------
% Standalone section: why the dichotomous goods result does not transfer to
% chores along the size shift. Uses only macros already defined in
% preamble.tex (\items, \agents, \alloc, \bundle, \cost, \subsidy,
% \optsubsidy, \edgew, \pathw, \set, \abs, \Z). Drop the file into
% sections/ and \input it wherever it is wanted.
% ---------------------------------------------------------------------------

\section{Why the dichotomous goods result does not transfer}
\label{sec:sizeshift}

The negative dichotomous class is carried onto the dichotomous class
of~\cite{BKNS22} by a single bijection, so it is natural to ask whether the
guarantee of that paper transfers directly. It does not. This section makes the
obstruction precise: the bijection preserves the valuation class but not the
envy graph, and the discrepancy is exactly the difference in bundle
cardinalities.

Throughout, $\cost_i : 2^{\items} \to \Z_{\ge 0}$ are dichotomous cost
functions, i.e.\ $\cost_i(\emptyset) = 0$ and
$\cost_i(S \cup \set{g}) - \cost_i(S) \in \set{0,1}$ for all $S$ and
$g \notin S$.

\subsection{The size shift}

\begin{definition}[Size shift]
\label{def:sizeshift}
For a cost function $\cost : 2^{\items} \to \Z_{\ge 0}$ define
$\widetilde{v} : 2^{\items} \to \Z$ by
\[
  \widetilde{v}(S) \;:=\; \abs{S} - \cost(S).
\]
\end{definition}

\begin{lemma}[The shift is a bijection of the two classes]
\label{lem:ss-bijection}
$\cost$ is a dichotomous cost function if and only if $\widetilde{v}$ is a
dichotomous valuation, i.e.\ $\widetilde{v}(\emptyset) = 0$ and every marginal
$\widetilde{v}(S \cup \set{g}) - \widetilde{v}(S)$ lies in $\set{0,1}$.
Moreover the map $\cost \mapsto \widetilde{v}$ is an involution.
\end{lemma}

\begin{proof}
$\widetilde{v}(\emptyset) = 0 - \cost(\emptyset) = 0$. For $g \notin S$,
\[
  \widetilde{v}(S \cup \set{g}) - \widetilde{v}(S)
  \;=\; \bigl(\abs{S}+1-\cost(S \cup \set{g})\bigr) - \bigl(\abs{S}-\cost(S)\bigr)
  \;=\; 1 - \bigl(\cost(S \cup \set{g}) - \cost(S)\bigr),
\]
so the marginal of $\widetilde{v}$ lies in $\set{0,1}$ exactly when that of
$\cost$ does. Applying the definition twice returns
$\abs{S} - (\abs{S} - \cost(S)) = \cost(S)$.
\end{proof}

\subsection{The shift does not preserve the envy graph}

For a complete allocation $\alloc = (\bundle1,\dots,\bundle n)$ write
$\edgew{\alloc}^{\,\cost}(i,j) = \cost_i(\bundle i) - \cost_i(\bundle j)$ for
the arc weights in cost form, and
$\edgew{\alloc}^{\,\widetilde{v}}(i,j) = \widetilde{v}_i(\bundle j) - \widetilde{v}_i(\bundle i)$
for those of the shifted instance.

\begin{lemma}[Arc-weight identity]
\label{lem:ss-arc}
For every allocation $\alloc$ and all $i,j \in \agents$,
\[
  \edgew{\alloc}^{\,\widetilde{v}}(i,j)
  \;=\; \bigl(\abs{\bundle j} - \abs{\bundle i}\bigr)
      \;+\; \edgew{\alloc}^{\,\cost}(i,j).
\]
\end{lemma}

\begin{proof}
By Definition~\ref{def:sizeshift},
\[
  \edgew{\alloc}^{\,\widetilde{v}}(i,j)
  = \widetilde{v}_i(\bundle j) - \widetilde{v}_i(\bundle i)
  = \bigl(\abs{\bundle j} - \cost_i(\bundle j)\bigr)
    - \bigl(\abs{\bundle i} - \cost_i(\bundle i)\bigr),
\]
and regrouping gives
$\bigl(\abs{\bundle j} - \abs{\bundle i}\bigr) + \bigl(\cost_i(\bundle i) - \cost_i(\bundle j)\bigr)$.
\end{proof}

\begin{corollary}
\label{cor:ss-equalcard}
The two envy graphs on $\alloc$ coincide arc by arc if and only if
$\abs{\bundle 1} = \dots = \abs{\bundle n}$. In that case $\alloc$ is
envy-freeable for $\cost$ if and only if it is envy-freeable for
$\widetilde{v}$, and the two minimum subsidy vectors agree.
\end{corollary}

\begin{proof}
Immediate from Lemma~\ref{lem:ss-arc}: the correction term
$\abs{\bundle j} - \abs{\bundle i}$ vanishes for all $i,j$ exactly when all
bundles have equal cardinality. Identical arc weights give identical cycles and
identical longest paths, and by
Theorems~\ref{thm:hs-characterisation} and~\ref{thm:hs-minsubsidy} both envy-freeability
and the minimum subsidy are read off the arc weights alone.
\end{proof}

Away from equal cardinalities the two minimum subsidies are related, but not
equal. Write $\widetilde{d}_{\alloc}(i,j)$ for the maximum weight of a directed
$i \to j$ path in the shifted envy graph.

\begin{proposition}
\label{prop:ss-subsidy}
Let $\alloc$ be envy-freeable for both instances. Then for every $i \in \agents$,
\[
  \pathw{\alloc}^{\,\cost}(i)
  \;=\; \max_{j \in \agents}\;
        \Bigl[\, \widetilde{d}_{\alloc}(i,j)
              + \abs{\bundle i} - \abs{\bundle j} \,\Bigr].
\]
Consequently, if $\widetilde{\subsidy}$ is any subsidy vector certifying
envy-freeness of $\alloc$ for $\widetilde{v}$ and we set
$q_i := \widetilde{\subsidy}_i + \abs{\bundle i}$, then
\[
  \optsubsidy_i \;=\; \pathw{\alloc}^{\,\cost}(i)
  \;\le\; q_i - \min_{j \in \agents} q_j .
\]
\end{proposition}

\begin{proof}
Sum the identity of Lemma~\ref{lem:ss-arc} along a directed path
$i = j_0 \to j_1 \to \dots \to j_k = j$. The correction terms telescope to
$\abs{\bundle{j}} - \abs{\bundle{i}}$, so the weight of the path in cost form
equals its weight in the shifted instance plus
$\abs{\bundle i} - \abs{\bundle j}$. Taking the maximum over $j$ and over paths
ending at $j$ gives the stated identity, since
$\pathw{\alloc}^{\,\cost}(i)$ is the maximum weight of a path leaving $i$ and
the empty path is the case $j=i$.

For the inequality, envy-freeness of $(\alloc,\widetilde{\subsidy})$ gives
$\widetilde{d}_{\alloc}(i,j) \le \widetilde{\subsidy}_i - \widetilde{\subsidy}_j$
for all $i,j$, whence
$\widetilde{d}_{\alloc}(i,j) + \abs{\bundle i} - \abs{\bundle j} \le q_i - q_j$.
Taking maxima over $j$ and applying
Theorem~\ref{thm:hs-minsubsidy} gives the claim.
\end{proof}

\subsection{What a transfer would require}

Proposition~\ref{prop:ss-subsidy} isolates the quantity that a transfer must
control. The guarantee of~\cite{BKNS22} produces, for the shifted instance, an
allocation with $\widetilde{\subsidy} \in \set{0,1}^n$; it asserts nothing about
$\abs{\bundle 1},\dots,\abs{\bundle n}$. The bound inherited on the chores side
is therefore governed not by $\widetilde{\subsidy}$ but by
$q_i = \widetilde{\subsidy}_i + \abs{\bundle i}$, and the per-agent conclusion
$\optsubsidy \in \set{0,1}^n$ follows only if $q$ has spread at most $1$.

\begin{proposition}
\label{prop:ss-target}
Suppose every dichotomous instance admits an allocation $\alloc$ and a
certifying subsidy vector $\widetilde{\subsidy}$ for which
$q_i = \widetilde{\subsidy}_i + \abs{\bundle i}$ satisfies
$\max_i q_i - \min_i q_i \le 1$. Then every negative dichotomous instance
admits an envy-free solution with subsidy in $\set{0,1}^n$.
\end{proposition}

\begin{proof}
Given costs $\cost_i$, apply the hypothesis to the shifted instance
$\widetilde{v}_i$, which is dichotomous by Lemma~\ref{lem:ss-bijection}, and
take the resulting $\alloc$. By Proposition~\ref{prop:ss-subsidy},
$\optsubsidy_i \le q_i - \min_j q_j \le 1$ for every $i$, and
$\optsubsidy_i \ge 0$ always.
\end{proof}

The hypothesis of Proposition~\ref{prop:ss-target} is strictly stronger than the
theorem of~\cite{BKNS22}, which constrains $\widetilde{\subsidy}$ alone, and it
is not supplied by that theorem: the allocation constructed there is not
guaranteed to be balanced, and indeed is not guaranteed to be EF1.

\subsection{The transfer fails on an explicit instance}

The following witness shows that the guarantee is not inherited allocation by
allocation, already for two agents and three chores with identical cost
functions.

\begin{example}
\label{ex:ss-fail}
Let $n = 2$, $\items = \set{a,b,c}$, and let both agents have
\[
  \cost_1(S) \;=\; \cost_2(S) \;=\; \max\bigl(0,\ \abs{S}-1\bigr).
\]
Every marginal is $0$ (from $\emptyset$ to a singleton) or $1$ (otherwise), so
this is a negative dichotomous instance. Its size shift is
\[
  \widetilde{v}_1(S) \;=\; \widetilde{v}_2(S)
  \;=\; \abs{S} - \max\bigl(0,\abs{S}-1\bigr) \;=\; \min\bigl(\abs{S},1\bigr),
\]
the unit-demand dichotomous valuation. Consider
\[
  \alloc \;=\; \bigl(\set{a,b,c},\ \emptyset\bigr).
\]

\emph{In the shifted goods instance.} We have
$\widetilde{v}_1(\bundle1) = 1$ and $\widetilde{v}_1(\bundle2) = 0$, so agent
$1$ does not envy agent $2$; and $\widetilde{v}_2(\bundle2) = 0$ with
$\widetilde{v}_2(\bundle1) = 1$, so agent $2$ envies agent $1$ by exactly $1$.
The allocation is envy-freeable and its minimum subsidy vector is
$\widetilde{\optsubsidy} = (0,1)$, consistent with the bound
of~\cite{BKNS22}.

\emph{In the chores instance.} We have $\cost_1(\bundle1) = 2$ and
$\cost_1(\bundle2) = 0$, so
\[
  \edgew{\alloc}^{\,\cost}(1,2) \;=\; \cost_1(\bundle1) - \cost_1(\bundle2) \;=\; 2 ,
\]
while $\edgew{\alloc}^{\,\cost}(2,1) = -2$. The allocation is envy-freeable, and
its minimum subsidy vector is $\optsubsidy = (2,0)$, which violates the
per-agent bound.
\end{example}

The two computations are reconciled by Lemma~\ref{lem:ss-arc}: here
$\abs{\bundle2} - \abs{\bundle1} = -3$, so
$\edgew{\alloc}^{\,\widetilde{v}}(1,2) = -3 + 2 = -1$, and in the notation of
Proposition~\ref{prop:ss-subsidy}, $q = (0 + 3,\ 1 + 0) = (3,1)$ has spread
$2$. The shifted instance is unaffected by the imbalance because the
correction absorbs it; the chores instance is not.

\begin{remark}
\label{rem:ss-scope}
Example~\ref{ex:ss-fail} does not contradict the main theorem. It shows that a
particular allocation, admissible on the goods side, is inadmissible on the
chores side, so the correspondence of Lemma~\ref{lem:ss-bijection} cannot by
itself carry the guarantee across; a separate argument is required to select the
allocation. By Corollary~\ref{cor:ss-equalcard} any such argument must control
bundle cardinalities, which is exactly the content that the theorem
of~\cite{BKNS22} does not provide.
\end{remark}
